\MWChapter{Mechanization and assurance}{机械验证与保障}
\label{ch:infrastructure}

\section{Formal project surface / 形式化项目范围}

The repository's \texttt{formal/} directory is a single Lean 4 proof project. Its lock file pins
Lean and mathlib, while the proof manifest connects central declarations to source integrity,
build evidence, and review metadata. The project gate rejects unknown kernel dependencies and
does not permit project theorems to rely on \texttt{sorry}, \texttt{admit}, project-defined axioms,
or \texttt{unsafe} declarations \citep{cantiluneFormalReadme}.

\begin{table}[H]
\centering
\caption{Mechanized assurance surfaces relevant to the theory report.}
\label{tab:infra}
\begin{tabularx}{\textwidth}{@{}L{0.25\textwidth}L{0.29\textwidth}Y@{}}
\toprule
\rowcolor{MWSoft!70}
\textbf{Surface} & \textbf{Recorded artifact} & \textbf{Assurance purpose} \\
\midrule
Theory source & Cantilune Lean source tree & Defines the mechanized categorical, rewriting, \(\pi\), feedback, and closure objects \\
Toolchain lock & Pinned Lean toolchain and resolved dependency lock & Makes the declared Lean and mathlib dependency state reconstructible \\
Proof manifest & Machine-readable central-obligation manifest & Names central obligations, Lean symbols, status, commits, and review fields \\
Source integrity & Source-integrity digest & Detects evidence/source mismatch before a proved claim is accepted \\
Build evidence & Durable QA build record & Preserves a durable audit trail rather than relying on a transient local cache \\
\bottomrule
\end{tabularx}
\end{table}

\section{Current recorded scale / 当前记录规模}

The branch README reports 565 maintained Lean source files, 1,624 audited declarations, and
18/18 central proof obligations recorded as \emph{proved}; these values describe the repository
evidence record, not a benchmark or a product reliability metric \citep{cantiluneReadme,cantiluneManifest}.

\begin{MWMetricGrid}
  \MWMetricCard[MWCyan]{565}{Maintained source files}{Lean source count in the branch record}
  \MWMetricCard[MWLilac]{1,624}{Audited declarations}{Recorded kernel-audit surface}
  \MWMetricCard[MWMint]{18 / 18}{Central obligations}{Manifest status: proved}
\end{MWMetricGrid}

\section{What the assurance gate means / 保障门的含义}

\begin{enumerate}
  \item \textbf{Kernel checking} establishes that the named Lean declarations elaborate under the
    pinned environment and the project gate's dependency restrictions.
  \item \textbf{Provenance binding} establishes that a status refers to a specific source and
    evidence record rather than an unpinned working tree.
  \item \textbf{Independent review} is not supplied by either of the first two steps. The manifest
    and delivery record retain explicit pending fields for reviewer and conflict-of-interest data.
  \item \textbf{Product conformance} is a later obligation: a concrete package must supply its
    own rule, projection, admission, resource, authorization, fairness, and replay certificates.
\end{enumerate}

\begin{MWCallout}[MWAmber]{No operational service claim}
This section concerns the integrity of formal artifacts. It is not infrastructure evidence for
availability, throughput, data protection, tool execution, or operational incident response,
because no such runtime is in scope for the current branch.
\end{MWCallout}
